\section{In-plane disorder leaves the transverse eigenvalue}
\label{sec:h-iso}

Proposition~\ref{prop:metric} is the local statement.
Random shifts and a transpose act on \((x,y)\).
The coefficient \(\ell_W^2\) of \(dz^2\) is invariant.
The sheet Hessian is already \(4I\), so those moves have no quadratic
in-plane anisotropy to remove, and they do not equalize the transverse
eigenvalue with the in-plane ones.

Proposition~\ref{prop:over} places the isotropic point \(\ell_W=1\) at
\(\lambda=1/4\), stub fraction \(1-e^{-1/2}\approx 0.393\), which the weave
preregistration calls overconnected.
Every woven value \(\lambda\le 0.08\) keeps \(\ell_W\ge 1.77\).
The volume exponent can equal \(3\) on that range
(Section~\ref{sec:h-weave}) while the axis ratio stays at least \(1.77\).

The average \(M_{ab}=\langle n_a n_b\rangle\) needs normals in an ambient
three-space.
The chain construction does not assign them.
You could read an isotropic infrared three-space, built from sheets, as a
meeting complex whose transverse and in-plane costs match.
On \eqref{eq:ds} that match is \(\ell_W=1\), which removes the scale
separation the crossover is built on.
